% 分类目录：dl
\documentclass[12pt, a4paper]{article}
\usepackage[margin=2.5cm]{geometry}
\usepackage{ctex}
\usepackage{amsmath, amssymb}
\usepackage{listings}
\usepackage{xcolor}
\usepackage{tabularx}
\usepackage{hyperref}

\lstset{
    language=Python,
    basicstyle=\ttfamily\small,
    keywordstyle=\color{blue},
    commentstyle=\color{green!50!black},
    stringstyle=\color{red},
    showstringspaces=false,
    breaklines=true,
    numbers=left,
    numberstyle=\tiny\color{gray},
    frame=single
}

\title{知识点：Qwen (通义千问)}
\author{}
\date{}

\begin{document}
\maketitle

\section{核心定义}
\textbf{中文定义}：Qwen（通义千问）是由阿里巴巴集团旗下的通义实验室（原达摩院）开发的大语言模型系列。其基于标准的 \textbf{Decoder-only Transformer} 架构，采用了多项现代优化技术，支持多语言、长文本以及极强的逻辑和编码能力。

\textbf{英文定义}：Qwen (Tongyi Qianwen) is a series of large language models developed by Alibaba Cloud's Tongyi Lab. Based on a standard Decoder-only Transformer architecture, it incorporates various modern optimization techniques and supports multilingual processing, long context, and strong reasoning and coding capabilities.

\textbf{核心特色}：
Qwen 强调在各项 Benchmark（如 MMLU, GSM8K, HumanEval）上的极致表现，目前通过开源（Qwen-7B/14B/72B）与闭源（Qwen-Max/Plus）并行，成为全球范围内最具竞争力的华语 LLM 之一。

\section{关键架构优化}
Qwen 的架构虽然遵循 Transformer Decoder，但进行了以下关键改进：

\begin{enumerate}
    \item \textbf{旋转位置编码 (RoPE)}：采用 Rotary Positional Embeddings 替代原始的绝对位置编码，增强了模型处理长序列的泛化能力。
    \item \textbf{SwiGLU 激活函数}：将标准的 FFN 激活函数从 ReLU/GeLU 升级为 SwiGLU，结合了 Gated Linear Unit 和 Swish 函数的优点，提升了模型的非线性理解力。
    \item \textbf{RMSNorm}：使用 Root Mean Square Layer Normalization 替代传统的 LayerNorm，只缩放均值而不偏移均值，计算效率更高且更稳定。
    \item \textbf{QKV Bias}：在多头注意力计算中加入了 Bias 项，进一步提高了模型的特征提取精度。
\end{enumerate}

\section{分词器 (Tokenizer)}
Qwen 采用了基于 \textbf{Tiktoken} 的分词器，词表规模高达 \textbf{151,851}。
\begin{itemize}
    \item 支持多国语言、代码片段以及特殊符号。
    \item 由于词表巨大，Qwen 在处理中文和复杂代码时具有极高的压缩效率（即同样的 token 能够承载更多的信息）。
\end{itemize}

\section{训练范式}
\begin{itemize}
    \item \textbf{Pre-training}：在数万亿 Token 的高质量中英文数据上进行自回归预训练。
    \item \textbf{SFT (Supervised Fine-tuning)}：利用高质量指令数据进行对齐。
    \item \textbf{RLHF (Alignment)}：通过人类反馈强化学习，特别是针对安全性、指令遵循和创意性进行细致打磨。
\end{itemize}

\section{Qwen 家族演进}
\begin{table}[h]
\centering
\begin{tabularx}{\textwidth}{|l|l|X|}
\hline
\textbf{版本} & \textbf{特色} & \textbf{应用场景} \\
\hline
Qwen-1.8B & 极小规模 & 移动端、端侧计算。 \\
\hline
Qwen-7B/14B & 主流开源 & 社群开发、个人研究。 \\
\hline
Qwen-72B & 最强开源 & 对标 GPT-3.5/Llama2-70B 的企业级应用。 \\
\hline
Qwen-VL & 多模态 & 图像描述、视觉文档分析。 \\
\hline
Qwen-Audio & 语音 & 语音识别与智能音频对话。 \\
\hline
\end{tabularx}
\end{table}

\section{代码演示：SwiGLU 概念实现}
\begin{lstlisting}
import torch
import torch.nn as nn
import torch.nn.functional as F

class SwiGLU(nn.Module):
    def __init__(self, d_model, d_ff):
        super().__init__()
        self.w1 = nn.Linear(d_model, d_ff)
        self.w2 = nn.Linear(d_model, d_ff)
        self.w3 = nn.Linear(d_ff, d_model)

    def forward(self, x):
        # Swish(x * w1) * (x * w2)
        swish_part = F.silu(self.w1(x))
        glu_part = self.w2(x)
        return self.w3(swish_part * glu_part)
\end{lstlisting}

\section{深度面试题}

\textbf{问：Qwen 为什么要采用 151,851 这么大的词表？}\\
答：大词表有三个主要优势：1. \textbf{编码效率}：同等长度的句子转化为更少的 Token，推理速度更快；2. \textbf{信息密度}：每个 Token 承载的语义信息更丰富；3. \textbf{兼容性}：更好地处理多种编程语言和罕见字符。

\textbf{问：RoPE（旋转位置编码）的核心数学原理是什么？}\\
答：RoPE 通过旋转特征向量的角度来注入位置信息，其最大的特点是\textbf{相对位置一致性}：两个位置的内积只取决于它们的相对距离，这使得模型在推理时能很容易地外推到比训练时更长的序列。

\section{参考文献}
\begin{enumerate}
    \item Bai, J., et al. (2023). \textit{Qwen Technical Report}. Alibaba Group.
    \item Touvron, H., et al. (2023). \textit{Llama 2: Open Foundation and Fine-Tuned Chat Models} (Qwen 架构与之相似).
    \item OpenAI. (2023). \textit{GPT-4 Technical Report}.
\end{enumerate}

\end{document}
